% Cumulative Telugu reader through OLP-0279.
% English source remains frozen under upstream/; editable Telugu remains under translation/.
\documentclass[11pt,a4paper,openany]{memoir}

\newcommand*{\olpath}{upstream}

\usepackage{fontspec}
\setmainfont{Latin Modern Roman}
\setsansfont{Latin Modern Sans}
\setmonofont{Latin Modern Mono}
\newfontfamily{\telugufont}{NotoSerifTelugu-Regular.ttf}[
  Path=fonts/,
  Script=Telugu,
  Language=Telugu,
  BoldFont=NotoSerifTelugu-Bold.ttf,
  ItalicFont=NotoSerifTelugu-Regular.ttf,
  BoldItalicFont=NotoSerifTelugu-Bold.ttf
]
\usepackage[Latin,Telugu]{ucharclasses}
\setTransitionsForLatin{\rmfamily}{}
\setTransitionTo{Telugu}{\telugufont}
\XeTeXgenerateactualtext=1

\input{\olpath/sty/open-logic.sty}
\input{\olpath/sty/open-logic-defer.sty}
\includeenv{editorial}
\problemsperchapter
\let\cleardoublepage\clearpage
\pagestyle{plain}

% Target-only evidence wrappers. The exact English token identity stays in the
% editable second argument but is deliberately absent from the reader surface.
\providecommand{\tetoken}[2]{#1}
% The canonical source writes the configurable negation token as `!\ `.
% OLP-0058 attaches the Telugu accusative suffix to that escaped-space token;
% defining the resulting control sequence preserves both pieces in print.
\providecommand{\ను}{{\telugufont ను}}
\providecommand{\MPని}{\MP{}\ను}
\providecommand{\QRతో}{\QR{}{\telugufont తో}}
\providecommand{\dotsను}{\dots{}\ను}
\providecommand{\dotsకు}{\dots{}{\telugufont కు}}
% Correction explanations remain verbatim in the editable TeX and evidence
% ledgers. The corrected reading is rendered, while the long audit annotation
% is omitted from running text so it remains safe inside conditionals/tables.
\newcommand{\sourcecorrection}[2]{}

\def\partname{{\telugufont భాగం}}
\def\chaptername{{\telugufont అధ్యాయం}}
\def\contentsname{{\telugufont విషయ సూచిక}}
\def\figurename{{\telugufont పటం}}
\def\proofname{{\telugufont నిరూపణ}}
\def\theoremname{{\telugufont సిద్ధాంతం}}
\def\lemmaname{{\telugufont ఉపసిద్ధాంతం}}
\def\propositionname{{\telugufont ప్రతిపాదన}}
\def\corollaryname{{\telugufont పర్యవసానం}}
\def\definitionname{{\telugufont నిర్వచనం}}
\def\examplename{{\telugufont ఉదాహరణ}}
\def\problemname{{\telugufont అభ్యాసం}}

\hypersetup{
  pdftitle={OpenLogic Telugu cumulative draft through OLP-0279},
  pdfauthor={Open Logic Project; Telugu machine translation},
  pdfsubject={276 of 722 editable source units; provisional terminology and review notes}
}
\tolerance=3000
\emergencystretch=4em
\raggedbottom

\newcounter{externalref}
\crefname{externalref}{{\telugufont పరిధికి వెలుపల సూచన}}{{\telugufont పరిధికి వెలుపల సూచనలు}}
\Crefname{externalref}{{\telugufont పరిధికి వెలుపల సూచన}}{{\telugufont పరిధికి వెలుపల సూచనలు}}

\begin{document}

\begin{titlingpage}
\begin{raggedleft}
{\HUGE\bfseries\sffamily\telugufont OpenLogic తెలుగు\par}
\vskip 3ex
{\huge\telugufont OLP-0279 వరకు సంచిత పాఠక ముసాయిదా\par}
\vskip 4ex
{\Large\telugufont 722 మూల విభాగాలలో 276 సంపాదించగల తెలుగు విభాగాలు\par}
\end{raggedleft}

\vfill

{\telugufont
ఇది పూర్తి గ్రంథం కాదు. అసంపూర్ణత భాగంలోని పరిచయ అధ్యాయం వరకు
అనువాదం కొనసాగింది; తరువాతి అధ్యాయాలు ఈ సంచికలో లేవు.

యంత్ర అనువాదం; మూల పాఠ్యంతో సరిపోల్చిన ఏజెంట్ సమీక్ష. మానవ లేదా
స్వతంత్ర నిపుణుల సమీక్ష జరిగిందని ఈ సంచిక పేర్కొనదు. సాంకేతిక పదాలు
తాత్కాలిక నిర్ణయాలు; ప్రదర్శిత నిర్వచనాలు, సూత్రాలే కచ్చితమైన అర్థాన్ని
నిర్ణయిస్తాయి.

ప్రతిజ్ఞావాక్యాత్మక, మొదటిస్థాయి సందర్భాలు పంచుకునే నిరూపణ-వ్యవస్థ
మాడ్యూళ్లను మూల గ్రంథపు సహజ దిగుమతి క్రమంలో రెండుచోట్ల ముద్రించవచ్చు.
}

\vskip 3ex
\small
Source revision: \texttt{9620cc73f9c8e0ad003c514a5d3748f29611c4c0}\par
Release ref: \texttt{v0.4.0-cumulative-olp0279}\par
\url{https://github.com/KokunoYumeto/OpenLogic-te-Telu-IN}
\end{titlingpage}

\frontmatter
\chapter*{{\telugufont సంచిక గురించి}}
{\telugufont
Open Logic Project మూలానికి Creative Commons Attribution 4.0 International
అనుమతి వర్తిస్తుంది. ఈ తెలుగు పాఠ్యం మూలంలో చేసిన మార్పు. మూల రచయితల
ఆమోదం ఉందని సూచించడం లేదు. 238 మూల దిద్దుబాటు వివరణలు సంపాదించగల
TeXలో, నిర్ణయ-సమీక్ష ఆధారాల్లో గుర్తింపులతో ఉన్నాయి; ఈ పాఠక PDFలో
సరిచేసిన పఠనం మాత్రమే కనిపిస్తుంది.

ఈ PDFలో 276 ప్రత్యేక సంపాదించగల TeX విభాగాల పరిధి ఉంది. వెబ్/EPUB
పాఠక సంచికకు వేరే, చిన్న, విఫలమైతే మూసివేసే మార్పిడి పరిధి ఉంది;
ఆ పరిధిని ఈ PDF సంఖ్యతో కలపకూడదు.
}

\section*{{\telugufont పరిధికి వెలుపల సూచనలు}}
{\telugufont కింది ఎనిమిది మూల సూచనలు ఈ PDFలో ఇంకా లేని భాగాలు లేదా
అధ్యాయాలకు వెళ్తాయి. ప్రశ్నార్థక సూచనలుగా చూపకుండా, వాటిని ఇక్కడ
స్పష్టంగా పరిధికి వెలుపలగా గుర్తించాం.}
\begin{enumerate}
\item\refstepcounter{externalref}\label{his:set:limits:sec}{\telugufont చారిత్రక సమితి-సిద్ధాంత పరిమితుల విభాగం.}
\item\refstepcounter{externalref}\label{his:set:mythology:sec}{\telugufont సమితి-సిద్ధాంత చరిత్ర/పురాణకథల విభాగం.}
\item\refstepcounter{externalref}\label{inc:req:min:lem:less-nsucc}{\telugufont అసంపూర్ణత భాగంలోని తరువాతి అవసరాల ఉపసిద్ధాంతం.}
\item\refstepcounter{externalref}\label{inc:req:min:lem:less-zero}{\telugufont అసంపూర్ణత భాగంలోని తరువాతి అవసరాల ఉపసిద్ధాంతం.}
\item\refstepcounter{externalref}\label{mth:ind:idf:sec}{\telugufont ఆగమనాత్మక నిర్వచనాల పద్ధతి విభాగం.}
\item\refstepcounter{externalref}\label{mth:ind:sti:sec}{\telugufont నిర్మాణాత్మక ఆగమన పద్ధతి విభాగం.}
\item\refstepcounter{externalref}\label{sth:::part}{\telugufont సమితి సిద్ధాంతం భాగం.}
\item\refstepcounter{externalref}\label{sth:ord-arithmetic::chap}{\telugufont ఆర్డినల్ అంకగణితం అధ్యాయం.}
\end{enumerate}

\tableofcontents*
\mainmatter

% OLP-0004--OLP-0054. The untranslated OLP-0003 part driver is outside
% the editable boundary, so its Telugu part heading is supplied here.
\olpart{sfr}{{\telugufont నైవ్ సమితి సిద్ధాంతం}}
\let\cumulativeoriginallabel\label
\newif\ifcumulativenatnatlabel
\renewcommand{\label}[1]{%
  \edef\cumulativelabelargument{#1}%
  \def\cumulativenatnatkey{sfr:siz:red:prob:nat-nat}%
  \ifx\cumulativelabelargument\cumulativenatnatkey
    \ifcumulativenatnatlabel
      \cumulativeoriginallabel{sfr:siz:red:prob:nat-nat-alt}%
    \else
      \global\cumulativenatnatlabeltrue
      \cumulativeoriginallabel{#1}%
    \fi
  \else
    \cumulativeoriginallabel{#1}%
  \fi}
\subfile{tmp/pdfs/cumulative-279/render-tree/content/sets-functions-relations/sets/sets}
\subfile{tmp/pdfs/cumulative-279/render-tree/content/sets-functions-relations/relations/relations-complete}
\subfile{tmp/pdfs/cumulative-279/render-tree/content/sets-functions-relations/functions/functions}
\subfile{tmp/pdfs/cumulative-279/render-tree/content/sets-functions-relations/size-of-sets/size-of-sets-complete}
\let\label\cumulativeoriginallabel
\subfile{tmp/pdfs/cumulative-279/render-tree/content/sets-functions-relations/arithmetization/arithmetization}
\subfile{tmp/pdfs/cumulative-279/render-tree/content/sets-functions-relations/infinite/infinite}

% Complete translated part drivers through the end of Undecidability.
\subfile{tmp/pdfs/cumulative-279/render-tree/content/propositional-logic/propositional-logic}
\subfile{tmp/pdfs/cumulative-279/render-tree/content/first-order-logic/first-order-logic}
\subfile{tmp/pdfs/cumulative-279/render-tree/content/model-theory/model-theory}
\subfile{tmp/pdfs/cumulative-279/render-tree/content/computability/computability}
\subfile{tmp/pdfs/cumulative-279/render-tree/content/turing-machines/turing-machines}

% OLP-0274 is retained, but its imports beyond the accepted boundary are
% suppressed. OLP-0275 then imports only the four translated intro sections.
\begingroup
\RenewDocumentCommand\olimport{s o m o}{}
\subfile{tmp/pdfs/cumulative-279/render-tree/content/incompleteness/incompleteness}
\endgroup
\subfile{tmp/pdfs/cumulative-279/render-tree/content/incompleteness/introduction/introduction}

\backmatter
\nocite{Frege1953}
\bibliographystyle{upstream/bib/natbib-oup}
\bibliography{upstream/bib/open-logic}

\end{document}
